% Laget automatisk av elevark/lag_elevark.py fra sim/berg-og-dal.html. Endringer her blir overskrevet.
\documentclass[11pt]{article}
\usepackage{elevark}
\begin{document}
\elevarkhode{Fysikk 1 · Mekanisk energi}{Berg-og-dal-bane}{https://fysikklab.no/sim/berg-og-dal.html}{Simuleringen}

Ei vogn får en liten dytt på toppen av en berg-og-dal-bane. Kommer den over de neste toppene og helt i mål? Den stiplede linja viser energinivået: så høyt kan vogna komme med den energien den har.

\section*{Oppgaver}

Jobb gjerne i par. Skriv ned gjetningen og begrunnelsen \emph{før} dere trykker Start.

\begin{oppgaveliste}
\oppgave \fase{gjett}{Gjett}Vogna starter 8,0~m over bakken med farten 1,0~m/s. Topp 2 er 7,5~m høy, og det er ingen friksjon. Kommer vogna over topp 2? Hvor stor er farten nede i den første dalen?
\svar
\oppgave \fase{test}{Test}Kjør simuleringen og sjekk tipsene. Se på grafen: Hva skjer med $E_{\mathrm{k}}$ når $E_{\mathrm{p}}$ øker?
\svar
\oppgave \fase{forklar}{Forklar}Hvor høy kan topp 2 være hvis vogna skal komme over? Regn ut svaret, og test med en topp litt lavere og litt høyere enn det dere fikk.
\svar
\oppgave \fase{test}{Regn}Regn ut farten vogna har når den kommer i mål. Endre høyden på topp 3. Endrer det farten i mål?
\svar
\oppgave \fase{gjett}{Gjett}Sett friksjonstallet til 0,01. Hvor på banen tror dere energinivået synker raskest? Test, og forklar hvorfor.
\svar
\oppgave \fase{utfordring}{Utfordring}Sett friksjonstallet til 0,02. Hvor høy kan topp 2 være nå hvis vogna skal komme i mål? Finn svaret ved å prøve dere fram, og forklar hvorfor det ikke er så lett å regne det ut.
\svar
\end{oppgaveliste}

\end{document}
